\chapter{Density operator}
%
\section{Density operator}
\subsection{Definition}
State vectors $\ket{\psi}$ convey the maximal amount of information about a system allowed by the laws of quantum mechanics.

If $\ket{\psi_1}$ and $\ket{\psi_2}$ are two possible quantum states then so is their coherent superporition
\begin{align*}
    \ket{\psi}=c_1\ket{\psi_1}+c_2\ket{\psi_2},\quad|c_1|^2+|c_2|^2=1.
\end{align*}
However, there are situations where the state vector is not precisely known.
It may be possible to write state vectors for the multicomponent system but not for the subsystem of interest.

For instance, a system of two spin-1/2 with eigenstates $\{\ket{\uparrow},\ket{\downarrow}\}$. A possible state of the combined system is 
\begin{align*}
    \ket{\psi}=\frac{1}{\sqrt{2}}\left[\ket{\uparrow}_1\ket{\downarrow}_2-\ket{\downarrow}_1\ket{\uparrow}_2\right].
\end{align*}
This equation is an entangled state. An entangled state cannot be factored, in any basis, into a product of states of the two subsystems, that is,
\begin{align*}
    \text{Entangled state}\qquad\ket{\psi}\neq\ket{\text{spin 1}}\ket{\text{spin 2}}.
\end{align*}
Quantum states described by state vectors are said to be \emph{pure states}. States that cannot be described by state vectors are said to be in 
\emph{mixed states}. Mixed states are described by the density operator 
\begin{align}
    \text{Mixed state}\qquad\highlight{\hrho=\sum_ip_i\ket{\psi_i}\bra{\psi_i}},
    \label{eq:densityoperator_mixed}
\end{align}
where the sym is over an ensemble and $p_i$ is the probability of the system being in the ith state of the ensemble $\ket{\psi_i}$.
The probabilities satisfy the obvious relations 
\begin{align*}
    0\leq p_i\leq1,\qquad\sum_ip_i=1,\qquad\sum_ip_i^2\leq1.
\end{align*}
For the pure state, all the $p_i$ vanish except the jth with $p_i\delta_{ij}$. Then,
\begin{align*}
    \text{Pure state}\qquad\highlight{\hrho=\ket{\psi_j}\bra{\psi_j}}.
    \label{eq:densityoperator_pure}
\end{align*}
The density operator for the pure state $\ket{\psi_j}$ is just the projection onto the state $\ket{\psi_j}$, while the general case 
\eqref{eq:densityoperator_mixed} is a sum of the projection operators over the ensemble, weighted with the probabilities of each member 
of the ensemble.

In an orthonormal basis $\{\ket{\varphi_n}\}$, the ith member of the ensemble can be projected as 
\begin{align*}
    \ket{\psi_i}=\sum_n\ket{\varphi_n}\braket{\varphi_n|\psi_i}=\sum_nc_n^{(i)}\ket{\varphi_n},\quad c_n^{(i)}=\braket{\varphi_n|\psi_i}.
\end{align*} 
The matrix element of $\hrho$ between the $n$ and $n'$ eigenstates is 
\begin{align*}
    \braket{\varphi_n|\hrho|\varphi_{n'}}=\sum_i\braket{\varphi_n|\psi_i}p_i\braket{\psi_i|\varphi_{n'}}=\sum_ip_ic_n^{(i)}c_{n'}^{(i)^*}.
\end{align*}
The quantities $\braket{\varphi_n|\hrho|\varphi_{n'}}$ form the elements of the density matrix. The trace of this matrix is 
\begin{align*}
    \Tr{\hrho}=\sum_n\braket{\varphi_n|\hrho|\varphi_n}=\sum_{i,n}\braket{\varphi_n|\psi_i}p_i\braket{\psi_i|\varphi_n}=\sum_{i,n}p_i\braket{\psi_i|\varphi_n}\braket{\varphi_n|\psi_i}=\sum_ip_i=1.
\end{align*}
Since $\hrho$ is Hermitian, the diagonal elements $0\leq\braket{\varphi_n|\hrho|\varphi_n}\leq1$.

Let us now consider the square $\hrho^2=\hrho\cdot\hrho$. For a pure state, we have 
\begin{align*}
    \hrho^2=\ket{\psi}\braket{\psi|\psi}\bra{\psi}=\ket{\psi}\bra{\psi}=\hrho\Longrightarrow\Tr{\hrho^2}=\Tr{\hrho}=1.
\end{align*}
For a statistical mixture 
\begin{align*}
    \hrho^2=\sum_{i,j}p_ip_j\ket{\psi_i}\braket{\psi_i|\psi_j}\bra{\psi_j}.
\end{align*}
Its trace is 
\begin{align*}
    \Tr{\hrho^2}=\sum_n\braket{\varphi_n|\hrho^2|\varphi_n}=\sum_{n,i,j}p_ip_j\braket{\varphi_n|\psi_i}\braket{\psi_i|\psi_j}\braket{\psi_j|\varphi_n}=\sum_{i,j}p_ip_j
    |\braket{\psi_i|\psi_j}|^2\leq\left[\sum_ip_i\right]^2-1.
\end{align*}
Thus, we have 
\begin{align}
    \Tr{\hrho^2}\begin{cases}
    =1,&\text{pure state}\\
    <1,&\text{mixed state}
    \end{cases}.
\end{align}

The ensemble average (or expectation value for pure state) of some operator $\hO$ is 
\begin{align}
    \text{Expectation}\qquad\highlight{\braket{\hO}=\Tr{\hrho\hO}=\begin{cases}
        \braket{\psi_i|\hO|\psi_i},&\text{pure state}\\
        \sum_ip_i\braket{\psi_i|\hO|\psi_i},&\text{mixed state}
    \end{cases}}.
\end{align}
Formally we can use the density operator to express the expectation as: 
\begin{align*}
    \braket{\hO}=\Tr{\hrho\hO}=\sum_n\braket{\varphi_n|\hrho\hO|\varphi_n}=\sum_{n,i}p_i\braket{\varphi_n|\psi_i}\braket{\psi_i|\hO|\varphi_n}=
    \sum_{i,n}p_i\braket{\psi_i|\hO|\varphi_n}\braket{\varphi_n|\psi_i}=\sum_ip_i\braket{\psi_i|\hO|\varphi_i}.
\end{align*}

%
\subsection{Entangled states}
Lets consider a two-particle system and that each particle can be in either of two one-particle states $\ket{\psi_1}$ or $\ket{\psi_2}$. Using the notation 
$\ket{\psi_i^{(j)}}$=particle $j$ in state $i$, we consider a pure two-particle superposition state 
\begin{align*}
    \ket{\psi}=C_1\ket{\psi_1^{(1)}}\otimes\ket{\psi_2^{(2)}}+C_2\ket{\psi_2^{(1)}}\otimes\ket{\psi_1^{(2)}}.
\end{align*}
This can be extended for multiparticle systems, where we can define the \emph{reduced density operators} for each of the subsystems by tracing the 
density operator over the states of all the other systems. The reduced density operator for particle 1 is 
\begin{align*}
    \hrho^{(1)}=\text{Tr}_2[\hrho]=\braket{\psi_1^{(2)}|\hrho|\psi_1^{(2)}}+\braket{\psi_2^{(2)}|\hrho|\psi_2^{(2)}}=|C_1|^2\ket{\psi_1^{(1)}}\bra{\psi_1^{(1)}}
    +|C_2|^2\ket{\psi_2^{(1)}}\bra{\psi_2^{(1)}}.
\end{align*}
This has the form of a mixed state for particle 1 as long as $C_i\neq0,\;i=1,2$. Simiarly, for particle 2,
\begin{align*}
    \hrho^{(2)}=\text{Tr}_1[\hrho]=|C_1|^2\ket{\psi_1^{(2)}}\bra{\psi_1^{(2)}}
    +|C_2|^2\ket{\psi_2^{(2)}}\bra{\psi_2^{(2)}}.
\end{align*}
When one of the particles is considered without regards to the other, it is generally in a mixed state. Thus one may characterized the degree of 
entanglement according to the degree of purity of either of the subsystems. 
\begin{itemize}[itemsep=0pt,topsep=0pt,parsep=0pt]
    \item $\Tr{\hrho^{(2)}}^2=1$, the state $\ket{\psi}$ is not an entangled state.
    \item $\Tr{\hrho^{(2)}}^2<1$, then $\ket{\psi}$ describes an entanglement between subsystems 1 and 2.
\end{itemize}





%%
\subsection{Dynamics of the density operator}
The density operator evolves unitarily according to 
\begin{align}
    \text{von Neumann eq 1}\qquad\partial_t\hrho=\frac{i}{\hbar}[\hrho,\hH],\quad\text{con}\quad i\hbar\partial_t\ket{\psi_i}=\hH\ket{\psi_i}.
\end{align}
Alternatively, we may write 
\begin{align*}
    \text{Heisenberg picture of $\hrho$}\qquad\hrho(t)=\hU(t,0)\hrho(0)\hU^\dagger(t,0),
\end{align*}
where the unitary evolution operator $\hU(t,0)$ satisfies the equation 
\begin{align}
    \qquad i\hbar\partial_t\hU=\hH\hU.
\end{align}
The equation is sometimes written as 
\begin{align}
    \text{von Neumann eq 2}\qquad\highlight{i\hbar\partial_t\hrho(t)=[\hH,\hrho(t)]=\mathcal{L}\hrho(t)},
\end{align}
where $\mathcal{L}$ is known as the Liouvilian superoperator.